\documentclass[11pt]{article}
\usepackage[a4paper,margin=28mm]{geometry}
\usepackage[T1]{fontenc}
\usepackage{lmodern,amsmath,amssymb,amsthm,mathtools,hyperref}
\newcommand{\ddc}{dd^c}
\newcommand{\PSH}{\operatorname{PSH}}
\newcommand{\MA}{\operatorname{MA}}
\newcommand{\PI}[2]{P_{#1}[#2]_{\mathcal I}}
\newtheorem{theorem}{Theorem}
\newtheorem{lemma}[theorem]{Lemma}
\newtheorem{proposition}[theorem]{Proposition}
\title{Transverse ideal recovery on the Hopf threefold}
\author{}\date{}
\begin{document}\maketitle
\noindent\textbf{Scope.} The criterion concerns the stated scalar Hopf
fibration and its full compatible systems. The unrestricted recovery
problem for arbitrary big backgrounds remains open.
\begin{abstract}
On the scalar Hopf threefold, we prove a recovery criterion using
positive initial lower singularity volume on the projective base.
The ambient manifold has unbounded upper Monge--Amp\`ere mass.
An elliptic-fiber supremum over all competitors gives an exact identity
between ambient and transverse ideal envelopes. The known recovery
theorem on the closed base then gives every finite slice and every
positive ambient perturbation. An infinite hyperplane series produces
a nonanalytic compatible system with a positive transverse lower
volume. Invariance of the positive scalar slices alone constructs the
full compatible transverse system; its existence is not a separate
assumption. We retain the distinction between closed transverse volumes
and nonclosed ambient masses in the modification and gauge statements.
\end{abstract}

Use $\ddc=i\partial\bar\partial/(2\pi)$ and
$d^c=i(\bar\partial-\partial)/(4\pi)$. Every potential has its genuine
canonical representative. The ideal $\mathcal I(tu)$ is defined by
local integrability of $|f|^2e^{-tu}$, at every real $t>0$.
Ideal equivalence means equality of all these ideals. We write
\[
 \PI{\theta}{u}
 =\left(\sup\{v\in\PSH(X,\theta):v\leq0,
                         v\sim_{\mathcal I}u\}\right)^*.
\]
The lower ordinary-type volume is the infimum of actual non-pluripolar
masses over every genuine potential ordinarily equivalent to the input.
It is not the mass of a selected representative in a nonclosed background.

\section{Geometry and the criterion}
Let
\[
 X=(\mathbf C^3\setminus\{0\})/\langle z\mapsto ez\rangle,
 \quad \pi:X\longrightarrow B=\mathbf P^2,
 \quad s=\log|z|^2.
\]
For the standard Fubini--Study form $\kappa$, normalized by
$\int_B\kappa^2=1$, put
\[
 \alpha=\pi^*\kappa=\ddc s,\qquad
 \beta=ds\wedge d^c s,\qquad h=\alpha+\beta.
\]
The derivatives of $s$ descend. The form $h$ is smooth positive,
$dh=-ds\wedge\alpha\ne0$, and the elliptic group
$E=\mathbf C/(\mathbf Z+2\pi i\mathbf Z)$ acts transitively on the
fibers and preserves $h$. Xia--Zhang \cite[Theorem 1.1]{XZ} construct
actual smooth strict potentials of divergent mass on this exact quotient.
Their operator and background equal $2\pi$ times ours, so dividing
their forms by $2\pi$ divides the masses by $(2\pi)^3$ and proves
$v_+(h)=\infty$. This genuine non-BMP example is an external dependency;
no all-dimensional extension is needed.

A full compatible system in background $\theta$ consists, for every
smooth positive Hermitian $\eta$, of normalized genuine ideal models
$\Gamma_\tau^\eta$ in $\theta+\eta$, $\tau<c$. The slices are decreasing
and pointwise concave in $\tau$, constant for $\tau\leq a<c$, and
\[
 \Gamma_\tau^\eta=\PI{\theta+\eta}{\Gamma_\tau^{\eta'}}
                 \quad\text{if }\eta\geq\eta'.
\]
Their decreasing limits in background $\theta$ are genuine ideal models
and independent of the auxiliary Hermitian metric. These are the original
system hypotheses. The known BMP recovery theorem states that positive
initial \emph{lower ordinary-type volume} forces recovery at every finite
slice. We use that theorem only on the closed projective base.

\begin{theorem}\label{thm:criterion}
Let $\Gamma$ be a full compatible system in background $h$ on $X$.
Suppose every positive scalar slice $\Gamma_\tau^{\epsilon h}$ is
invariant under $E$, for all $\epsilon>0$ and $\tau<c$.
Let $\psi_\tau$ be the canonical base descent of the genuine decreasing
scalar limit $\Gamma_\tau$. If $\psi_a$ has positive lower
ordinary-type volume in $\kappa$, then, for every finite $\tau<c$
and every positive Hermitian $\eta$ on $X$,
\[
 \Gamma_\tau=\pi^*\psi_\tau,\qquad
 \Gamma_\tau^\eta=\PI{h+\eta}{\Gamma_\tau}.
\]
The positive scalar slices canonically determine a full compatible system
in background $\kappa$, for every positive Hermitian base perturbation.
No equality of their ideals with the limiting ideals is assumed.
\end{theorem}

\section{Local ideal constraints and envelope closure}
The local facts used below also apply after bounded smooth corrections.
In Boucksom's convention \cite{Bou}, $\mathcal J(\phi)$ uses
$|f|^2e^{-2\phi}$, so $\mathcal I(tu)=\mathcal J(tu/2)$.
For every local prime divisor $F$ over a smooth chart,
write $\nu_F$ in the function Lelong convention
$\nu_F(\log|z|)=1$. In particular
$\nu_F(\log|f|^2)=2\operatorname{ord}_F(f)$; the coefficient of
$[F]$ in our $\ddc\phi$ is $\nu_F(\phi)/2$. With this convention,
\begin{equation}\label{eq:orders}
 \nu_F(\phi)-\frac{A(F)}k
 \leq\frac1k\operatorname{ord}_F\mathcal J(k\phi)
 \leq\nu_F(\phi).
\end{equation}
The upper bound is \cite[Theorem 9.17]{Bou}: the analytic potential of
$\mathcal J(k\phi)^{1/k}$ is less singular than $\phi$. For the lower
bound, change variables near a regular point of $F$ and apply
\cite[Lemma 10.17 and Theorem 10.13]{Bou} to each generator; the Jacobian
has order $A(F)-1$. Coherence permits finitely many generators.
Thus equality of all multiplier ideals implies equality of all local
divisorial orders by letting $k\to\infty$. The converse is
\cite[Corollary 10.18]{Bou}, with its strong openness qualification.

Lelong numbers are upper semicontinuous under local $L^1$ convergence
\cite[Lemma 2.6]{Bou}. Their generic value on a divisor is the infimum
of the pointwise values, and
\[
 \nu_F(\max(u,v))=\min(\nu_F(u),\nu_F(v)),\qquad
 \nu_F(tu+(1-t)v)=t\nu_F(u)+(1-t)\nu_F(v)
\]
by \cite[Corollary 2.10 and Section 10.2]{Bou}.
For a holomorphic submersion $f$ we also have
\begin{equation}\label{eq:submersion}
 \mathcal I(tf^*u)=f^*\mathcal I(tu)\mathcal O.
\end{equation}
In local product Stein disks, Demailly's full Hilbert-tensor-product
argument \cite[Propositions 5.7 and 14.2(i)]{Dem} with fiber weight zero
proves this for all holomorphic germs. The reverse Fubini test on lifted
germs also shows that ideal equality upstairs implies equality downstairs.

\begin{lemma}\label{lem:closure}
Let $u\leq0$ be genuine $\theta$-psh. The nonempty ideal envelope
$\PI{\theta}{u}$ retains every ideal of $u$, is a genuine normalized
ideal model, and satisfies the nested-background identity
\[
 \PI{\theta_2}{\PI{\theta_1}{u}}
          =\PI{\theta_2}{u}\quad(\theta_1\leq\theta_2).
\]
\end{lemma}
\begin{proof}
Choquet selection realizes the envelope as the increasing canonical limit
of finite maxima, including a fixed competitor. Finite maxima preserve
the smooth-background lower Hessian: local convolution and a smooth convex
maximum give the same bound with an error tending to zero. Distributional
passage retains the exact bound. A fixed integrable lower competitor
prevents an identically minus-infinite limit.

Pull the sequence back to any local modification. Its increasing envelope
equals the pullback of the ambient envelope on the isomorphism locus,
then everywhere by canonical extension. Common smooth corrections give
local $L^1$ convergence there. Every common generic divisorial lower bound
survives by pointwise upper semicontinuity. The limit dominates a fixed
competitor, giving the reverse order inequality. Equation
\eqref{eq:orders} and the valuative converse prove exact ideal equivalence.
The envelope is its own model because its competitor family is unchanged.
Its supremum is zero: otherwise adding the negative of that supremum
would be another larger competitor. Finally the two inputs in the nested
identity have identical ideals, hence identical competitor families in
$\theta_2$. No mass hypothesis enters this argument.
\end{proof}

\section{The full fiber-envelope identity}
\begin{proposition}\label{prop:fiber}
For $b>0$ and every normalized genuine $b\kappa$-psh $v$ on $B$,
\[
 \PI{bh}{\pi^*v}=\pi^*\PI{b\kappa}{v}.
\]
\end{proposition}
\begin{proof}
Take any ambient competitor $q$. Every elliptic translate $q_g$ has
the same background and the same ideals as $\pi^*v$, by
\eqref{eq:submersion}. Lemma \ref{lem:closure} applied to the orbit
family shows that $M=(\sup_gq_g)^*$ retains those ideals. It is invariant
under every fiber translation and descends to an usc function $w$.

This descent gives the exact current identity
\[
 bh+\ddc M=\pi^*(b\kappa+\ddc w)+b\beta.
\]
In a holomorphic product chart lift any base vector smoothly to
$\ker(\partial s)$. The $\beta$ pairing is zero. Test positivity with
a positive smooth fiber density of integral one; since current coefficients
are measures, this proves the distributional positive pairing of
$b\kappa+\ddc w$ with that base vector. Fubini gives local $L^1$
integrability. A local potential of $b\kappa$ gives the canonical base
psh representative. Its pullback and $M$ have the same distribution;
canonical uniqueness identifies them everywhere, including poles.
Thus $w$ is genuinely $b\kappa$-psh.

Equation \eqref{eq:submersion} makes $w$ ideal-equivalent to $v$, so
$q\leq M=\pi^*w\leq\pi^*\PI{b\kappa}{v}$. Taking all ambient
competitors proves one inequality. Every base competitor lifts by the
displayed positive-current decomposition, giving the other.
\end{proof}

\section{Construction from scalar invariance}
\begin{lemma}\label{lem:scalar}
Under the scalar invariance hypothesis of Theorem \ref{thm:criterion},
there is a canonical full compatible base system $\Psi$ with
\[
 \Gamma_\tau^{\epsilon h}=\pi^*\Psi_\tau^{\epsilon\kappa},
 \qquad \Gamma_\tau=\pi^*\psi_\tau.
\]
Its genuine decreasing limit is independent of the auxiliary base metric.
\end{lemma}
\begin{proof}
The horizontal positive-current test in Proposition \ref{prop:fiber}
descends each scalar slice as
$\Gamma_\tau^{\epsilon h}=\pi^*\phi_\tau^\epsilon$, with
$\phi_\tau^\epsilon$ genuinely $(1+\epsilon)\kappa$-psh, including
its canonical values at poles. Normalization, decrease, pointwise concavity
and constancy for $\tau\leq a$ descend. The ambient model identity and
Proposition \ref{prop:fiber} show that each $\phi_\tau^\epsilon$ is
its own base ideal model. Ambient compatibility gives
\[
 \phi_\tau^{\epsilon_2}
 =\PI{(1+\epsilon_2)\kappa}{\phi_\tau^{\epsilon_1}}
 \quad(\epsilon_2\geq\epsilon_1>0).
\]
Thus Lemma \ref{lem:closure} gives equality of every positive-weight
ideal between these positive scalar slices. This concerns the positive
perturbations, not the decreasing limit.

For a positive smooth Hermitian base form $\xi$, choose $\delta>0$ with
$\delta\kappa\leq\xi$, fixed for the whole $\tau$ interval, and set
\[
 \Psi_\tau^\xi=\PI{\kappa+\xi}{\phi_\tau^\delta}.
\]
The input is admissible, so the envelope is nonempty. Lemma
\ref{lem:closure} makes it a normalized genuine ideal model with precisely
the input ideals. The preceding scalar identity shows that any other
allowed $\delta$ has the same ideals, hence the same complete competitor
family. The definition is independent of $\delta$. It reproduces
$\Psi_\tau^{\epsilon\kappa}=\phi_\tau^\epsilon$, and equality of
the ideal classes gives
\[
 \PI{\kappa+\xi}{\Psi_\tau^{\xi'}}=\Psi_\tau^\xi
 \quad(\xi\geq\xi').
\]
Neither argument requires the perturbing form $\xi$ to be closed.

Fix $\xi$ and an allowed $\delta$. If $\tau_1\leq\tau_2$, every
later competitor $v$ has the later target divisorial orders. The earlier
input $\phi_{\tau_1}^\delta$ has no larger order at any local divisor
over any chart. Consequently
$\max(v,\phi_{\tau_1}^\delta)$ is an admissible earlier competitor:
its order is exactly the earlier target at every such divisor. The
valuative characterization proves equality of all its ideals. Taking
the complete suprema proves decrease in $\tau$.

For $\tau=(1-t)\tau_1+t\tau_2$ and $0<t<1$, scalar concavity gives
\[
 \nu_F(\phi_\tau^\delta)
 \leq(1-t)\nu_F(\phi_{\tau_1}^\delta)
       +t\nu_F(\phi_{\tau_2}^\delta)
\]
at every local divisor $F$. For any endpoint competitors $v_1,v_2$, the
function $\max((1-t)v_1+tv_2,\phi_\tau^\delta)$ is admissible and
has exactly the intermediate target orders. It is therefore an
intermediate ideal-equivalent competitor dominating their convex sum.
Taking both full suprema proves pointwise concavity of $\Psi$.
Choquet selection, finite maxima and local canonical extension justify
these inequalities after regularization, including poles. At $t=0,1$
omit the term with zero coefficient. Constancy on the left is inherited
from the scalar inputs.

The original system supplies a genuine decreasing limit $\Gamma_\tau$.
It remains invariant. The same horizontal test descends it as
$\pi^*\psi_\tau$, with $\psi_\tau$ genuinely $\kappa$-psh.
The original ambient limit model identity and Proposition
\ref{prop:fiber} at $b=1$ make $\psi_\tau$ its own base ideal model.
Normalization also follows from nested nonempty compact superlevel sets.
For any positive smooth base metric $k$, choose $0<l\leq r$ with
$l\kappa\leq k\leq r\kappa$. The established compatibility and
envelope monotonicity yield
\[
 \phi_\tau^{l/j}\leq\Psi_\tau^{k/j}\leq\phi_\tau^{r/j}.
\]
The middle sequence decreases; both scalar sequences decrease to
$\psi_\tau$. This proves the genuine limit and independence of $k$.
No survival of the positive-slice ideals in that limit was used.
\end{proof}

\begin{proof}[Proof of Theorem \ref{thm:criterion}]
Lemma \ref{lem:scalar} constructs the full compatible base system from
scalar invariance alone.
The task's known positive-initial-lower-volume BMP theorem applies to
the full base system on $B$, yielding
\[
 \Psi_\tau^{\epsilon\kappa}
       =\PI{(1+\epsilon)\kappa}{\psi_\tau}
\]
for every finite $\tau<c$. Proposition \ref{prop:fiber} gives recovery
at every scalar ambient perturbation and identifies the canonical scalar
limit with $\pi^*\psi_\tau$. For arbitrary positive $\eta$, choose
$\epsilon h\leq\eta$. Compatibility and Lemma \ref{lem:closure} give
\[
 \Gamma_\tau^\eta
 =\PI{h+\eta}{\Gamma_\tau^{\epsilon h}}
 =\PI{h+\eta}{\PI{(1+\epsilon)h}{\Gamma_\tau}}
 =\PI{h+\eta}{\Gamma_\tau}.
\]
This uses no finite upper-volume deficit on $X$.
\end{proof}

\section{A nonanalytic compatible example}
Choose infinitely many distinct hyperplanes $H_j$ on $\mathbf P^2$,
$q_j>0$ with $\sum_jq_j=1$, and normalized standard logarithms
$\ell_j\leq0$ with $\kappa+\ddc\ell_j=[H_j]$. Unitary invariance
bounds their $L^1$ norms uniformly. The negative series
\[
 \psi_0=\sum_jq_j\ell_j-\sup_B\sum_jq_j\ell_j
\]
is a nontrivial genuine $\kappa$-psh function and
$\kappa+\ddc\psi_0=\sum_jq_j[H_j]$.
Fix $0<\lambda_0<\lambda_{\max}<1$ and choose $r>0$ so that
\[
 \lambda(\tau)=\lambda_0+r\max(\tau-a,0)<\lambda_{\max}
                      \quad(\tau<c),\qquad u_\tau=\lambda(\tau)\psi_0.
\]
Then $\kappa+\ddc u_\tau\geq(1-\lambda_{\max})\kappa$.
Positive orders along infinitely many distinct prime divisors preclude
an analytic singularity type, since a coherent ideal on a compact
manifold has finitely many divisorial components.

Define the full families, for every positive $\xi$ and $\eta$, by
\[
 \Psi_\tau^\xi=\PI{\kappa+\xi}{u_\tau},\qquad
 \Gamma_\tau^\eta=\PI{h+\eta}{\pi^*u_\tau}.
\]
The strict input currents make the families nonempty. Lemma
\ref{lem:closure} gives normalized genuine ideal models and full
compatibility; Proposition \ref{prop:fiber} gives the scalar basic identity.

For monotonicity, if $\tau_1\leq\tau_2$, the maximum of any later
competitor and $u_{\tau_1}$ has precisely the earlier target orders and
dominates that competitor. The same proof holds upstairs. For concavity,
combine two competitors in one background. Its order is the corresponding
affine combination of $\lambda(\tau_i)\nu_F(\psi_0)$. Convexity of
$\lambda$ makes this at least the intermediate target order. Taking the
maximum with the actual intermediate input restores equality of every
order and dominates the combination. Taking the full suprema proves
pointwise concavity. This argument retains all exceptional and
accumulation-point valuations, not just the listed hyperplanes.

The scalar base envelopes decrease as $\epsilon\downarrow0$ and are
bounded below by $u_\tau$. Their genuine limit has orders at least the
target by local $L^1$ semicontinuity and at most the target by that lower
bound. It is therefore $\PI{\kappa}{u_\tau}$: it is a competitor,
and each scalar envelope dominates this smaller-background envelope.
Supremum zero passes to the limit by compactness and decreasing closed
superlevel sets. Scalar ambient limits follow from the fiber identity;
squeezing arbitrary perturbing metrics between scalar multiples of $h$
proves background independence. The limits are genuine ideal models.

For the initial base model $\psi_a=\PI{\kappa}{\lambda_0\psi_0}$,
positivity and finite-level non-pluripolar multilinearity give
\[
 \int_B\MA_\kappa(\lambda_0\psi_0)\geq(1-\lambda_0)^2.
\]
The pole set has zero smooth volume. Since $\psi_a\geq\lambda_0\psi_0$,
closed-background mass monotonicity \cite[Theorem 1.2]{WN} gives the
same lower bound for $\psi_a$. Applied in both directions to every
ordinarily equivalent competitor, it proves
\[
 \underline{\operatorname{vol}}(\kappa,\psi_a)
 =\int_B\MA_\kappa(\psi_a)\geq(1-\lambda_0)^2>0.
\]
Thus this is a lower \emph{type} volume, not a selected ambient mass.
The example has nonanalytic ideal data on a genuinely unbounded-upper-mass
ambient manifold and meets the stated transverse criterion.

\section{Modification and gauge scope}
On a compact K\"ahler modification of the projective base, actual
non-pluripolar products pull back on the isomorphism locus and give no
mass to the exceptional analytic set. Potentials in the pulled background
descend: their restrictions to compact connected fibers are psh and
constant, possibly identically minus infinity; descent off the exceptional
locus and canonical extension identify the representatives. Ordinary
competitor types therefore correspond and the closed transverse lower
volume is preserved. Ideal-envelope compatibility is
\cite[Lemma 3.1]{DX}, with the exact pulled closed background; its
change-of-variables step is \cite[Proposition 5.8]{Dem}.

For arbitrary ambient modifications the conclusion preserves every local
divisorial order, by the local ideal characterization above. We do not
identify nonclosed numerical volumes on those modifications or assert
recovery for arbitrary new perturbation systems there. The modified
manifold need not be another scalar Hopf fibration.

Under the original smooth gauge transport, the criterion is transported
with its geometric data and normalization; the known gauge theorem
preserves the proved recovery. A newly gauged normalized system need not
have basic scalar slices. The classical BMP branch remains its own known
positive-initial-ambient-lower-volume theorem. These statements retain
the scope of the new transverse branch.

\begin{thebibliography}{9}
\bibitem{Bou} S. Boucksom,
\emph{Singularities of plurisubharmonic functions and multiplier ideals},
course notes, Sections 2, 9 and 10.
\bibitem{Dem} J.-P. Demailly,
\emph{Multiplier ideal sheaves and analytic methods in algebraic geometry},
Trieste lecture notes, Propositions 5.7--5.8 and 14.1--14.3.
\bibitem{WN} D. Witt Nystr\"om,
\emph{Monotonicity of non-pluripolar Monge--Amp\`ere masses},
arXiv:1703.01950v1, Theorem 1.2.
\bibitem{DX} T. Darvas and M. Xia,
\emph{The volume of pseudoeffective line bundles and partial equilibrium},
Geometry \& Topology 28 (2024), 1957--1993, Lemma 3.1.
\bibitem{XZ} M. Xia and K. Zhang,
\emph{A counterexample to the bounded mass property},
arXiv:2608.21053v1.
\end{thebibliography}
\end{document}
